\documentclass[a4paper,10pt]{article}
\usepackage[utf8]{inputenc}
\usepackage[french]{babel}
\usepackage[margin=1cm]{geometry}
\usepackage{multicol}
\usepackage{tikz}
\usepackage{xcolor}
\usepackage[T1]{fontenc}

\definecolor{parasitecolor}{RGB}{120,150,100}
\definecolor{safecolor}{RGB}{80,120,150}

\newcommand{\ficheparasite}[1]{
\colorbox{parasitecolor}{\textcolor{white}{\textbf{#1}}}
}

\newcommand{\fichenonparasite}[1]{
\colorbox{safecolor}{\textcolor{white}{\textbf{#1}}}
}

\pagestyle{empty}
\setlength{\parindent}{0pt}
\setlength{\parskip}{3pt}

\begin{document}

% ==================== PAGE 1 : JOUR 3 ====================
\noindent
\begin{minipage}[t]{0.48\textwidth}
\raggedright
\ficheparasite{JOUR 3 - TU ES PARASITÉ·E}

\vspace{0.3cm}

Pendant la nuit, quelque chose a changé.

\vspace{0.2cm}

Un réseau fongique s'est délicatement connecté à ton système nerveux. Tu ne ressens aucune douleur. Au contraire : une sérénité profonde t'envahit, comme si un poids invisible venait de quitter tes épaules.

\vspace{0.2cm}

\textbf{Cette planète est ton foyer maintenant.}

\vspace{0.2cm}

Les autres parlent encore de réparations, de décollage, de retour. Ces mots te semblent étrangement... creux. Pourquoi partir ? Vous avez enfin trouvé un endroit où vous pourriez être heureux. Où vous pourriez appartenir à quelque chose de plus grand.

\vspace{0.2cm}

\textbf{Tes nouveaux objectifs :}

\begin{itemize}
\setlength\itemsep{0.1em}
\item \textbf{Empêcher le départ}.
\item \textbf{Contaminer les autres}.
\subitem {Réussis à soigner une personne à l'issue d'un jet de dés particulièrement favorable. Débrouille toi pour trouver de quoi soigner.}
\subitem {Lors des actions collectives cruciales, essaie de rater complètement tes jets de dés.}
\item \textbf{Reste naturel} : ne révèle jamais ouvertement ta condition.
\end{itemize}

\vspace{0.2cm}

\textit{Tu n'es pas méchant. Tu crois sincèrement que rester est la meilleure chose pour tout le monde.}

\vspace{0.2cm}

Le mycélium murmure doucement : \textbf{« Bienvenue à la maison. »}
\end{minipage}
\hfill
\begin{minipage}[t]{0.48\textwidth}
\raggedright
\fichenonparasite{JOUR 3 - TOUT VA BIEN...}

\vspace{0.3cm}

Tu t'es réveillé·e normalement ce matin. Le soleil étrange de cette planète filtre à travers l'atmosphère. Tout semble en ordre.

\vspace{0.2cm}

Mais quelque chose a changé.

\vspace{0.2cm}

C'est subtil. À peine perceptible. Tu sens un nouveau regard peser dans le creux de ton cou. Certaines phrases sonnent... faux ? Ou est-ce ton imagination ?

\vspace{0.2cm}

\textbf{Tu n'es pas paranoïaque. Tu es prudent·e.}

\vspace{0.2cm}

Peut-être que tout va bien. Peut-être que c'est la fatigue du crash, le stress qui retombe.

\vspace{0.2cm}

Ou peut-être pas.

\vspace{0.2cm}

\textbf{Reste vigilant.}

\vspace{0.2cm}

\textbf{Ne baisse pas ta garde.}
\end{minipage}

\newpage

% ==================== PAGE 2 : JOUR 4 ====================
\noindent
\begin{minipage}[t]{0.48\textwidth}
\raggedright
\ficheparasite{JOUR 4 - TU ES PARASITÉ·E}

\vspace{0.3cm}

\textbf{Tes nouveaux objectifs :}

\begin{itemize}
\setlength\itemsep{0.1em}
\item \textbf{Empêcher le départ}.
\item \textbf{Contaminer les autres}.
\subitem {Réussis à soigner une personne à l'issue d'un jet de dés particulièrement favorable. Débrouille toi pour trouver de quoi soigner.}
\subitem {Lors des actions collectives cruciales, essaie de rater complètement tes jets de dés.}
\item \textbf{Reste naturel} : ne révèle jamais ouvertement ta condition.
\end{itemize}

\vspace{0.2cm}

Le réseau fongique pulse plus fort maintenant. Tu sens son urgence.

\vspace{0.2cm}

\textbf{Le temps presse.}

\vspace{0.2cm}

Tu as remarqué les préparatifs. Ils veulent partir. Ils parlent de réparations, de cartes, de trajectoires de fuite. Chaque discussion technique te serre le cœur. Ne comprennent-ils pas qu'ils ont déjà trouvé le paradis ?

\vspace{0.2cm}

\textbf{Ne perds pas de vue tes objectifs}

\vspace{0.2cm}

\textit{Il ne reste peut-être que deux jours pour les sauver malgré eux.}

\vspace{0.2cm}

Le mycélium murmure : \textbf{« Bientôt, vous serez tous chez vous. »}
\end{minipage}
\hfill
\begin{minipage}[t]{0.48\textwidth}
\raggedright
\fichenonparasite{JOUR 4 - RESTE VIGILANT}

\vspace{0.3cm}

Ce n'était pas ton imagination.

\vspace{0.2cm}

Les signes se multiplient. Des gestes inhabituels. Des hésitations inexplicables lors d'actions critiques.

\vspace{0.2cm}

\textbf{Qui sont-ils ?}

\vspace{0.2cm}

Tu ne peux accuser sans preuve. Mais tu ne peux non plus faire confiance aveuglément.

\vspace{0.2cm}

\textbf{Reste lucide.}

\vspace{0.2cm}

\textit{Chaque heure compte. Le vaisseau doit partir. Mais comment fuir quand l'ennemi dort dans la couchette d'à côté ?}

\vspace{0.3cm}

\textbf{Fais confiance à ton instinct.}
\end{minipage}

\newpage

% ==================== PAGE 3 : JOUR 5 ====================
\noindent
\begin{minipage}[t]{0.48\textwidth}
\raggedright
\ficheparasite{JOUR 5 - C'EST MAINTENANT OU JAMAIS}

\vspace{0.2cm}

Tu peux désormais te révéler et engager frontalement le combat contre les autres PJ (personnages joueurs) si tu le souhaites.

\vspace{0.3cm}

{\Large \textbf{AUJOURD'HUI.}}

\vspace{0.2cm}

Le vaisseau pourrait décoller demain. Peut-être même ce soir.

\vspace{0.2cm}

Le réseau fongique vibre d'une énergie désespérée. Tu ne peux plus te permettre la subtilité.

\vspace{0.2cm}

\textbf{TOUT DONNER. MAINTENANT.}

\vspace{0.2cm}

\textbf{Rappel de tes objectifs :}
\begin{itemize}
\setlength\itemsep{0.1em}
\item \textbf{Contaminer les autres} (toujours prioritaire) :
\subitem {Soigne quelqu'un à l'issue d'un jet de dés particulièrement favorable.}
\subitem {Rate complètement tes jets lors des actions collectives cruciales.}
\item \textbf{Empêcher le départ} (encore plus urgent).
\end{itemize}

\vspace{0.2cm}

\textbf{Ta mission finale :}

\begin{itemize}
\setlength\itemsep{0.1em}
\item \textbf{Sabotage majeur} : le vaisseau ne doit PAS pouvoir partir. Sois créatif pour empêcher le décollage.
\item \textbf{Rallie-toi aux autres parasités} : tu n'es plus seul. Coordonnez-vous. Des regards suffisent.
\item \textbf{Joue un rôle} : tu dois choisir entre ces deux possibilités role play (obligatoire) :
  \begin{itemize}
  \item Trouve toi un TIC comportemental obssessionnel.
  \item Une phrase un peu récurrente que tu ressors trop souvent.
  \end{itemize}
\end{itemize}

\vspace{0.2cm}

\textit{Si le vaisseau décolle, et que tu es dedans, tu mourras.}

\vspace{0.2cm}

{\Large \textbf{Ne les laisse pas partir.}}
\end{minipage}
\hfill
\begin{minipage}[t]{0.48\textwidth}
\raggedright
\fichenonparasite{JOUR 5 - MÉFIANCE MAXIMUM}

\vspace{0.2cm}

{\Large \textbf{C'EST AUJOURD'HUI QU'ON PART.}}

\vspace{0.2cm}

Ou jamais.

\vspace{0.2cm}

Les comportements étranges ont atteint un pic. Certains membres de l'équipage ne se cachent presque plus. Ils parlent avec une sérénité dérangeante. Ils sourient trop. Ils sabotent.

\vspace{0.2cm}

\textbf{C'est une lutte pour la survie.}

\vspace{0.2cm}

\textbf{Protocole d'urgence :}

\begin{itemize}
\setlength\itemsep{0.1em}
\item \textbf{Ne fais confiance à PERSONNE} : même tes alliés d'hier peuvent avoir changé cette nuit.
\item \textbf{Reste en groupe (sain)} : si tu identifies d'autres personnes lucides, collez-vous. Le nombre fait la force.
\item \textbf{Bats-toi si nécessaire} : si quelqu'un agit dangereusement, tu n'as plus le luxe du doute poli.
\end{itemize}

\vspace{0.2cm}

\textit{Cette planète veut vous garder. Elle a déjà pris certains d'entre vous.}

\vspace{0.2cm}

\textbf{Ne devient pas le prochain.}

\vspace{0.2cm}

\textit{Chaque geste compte. Chaque seconde compte.}

\vspace{0.2cm}

{\Large \textbf{Rentre chez toi. Vivant.}}
\end{minipage}

\end{document}
